% Part: proof-theory
% Chapter: cut-elimination
% Section: interpolation

\documentclass[../../../include/open-logic-section]{subfiles}

\begin{document}

\olfileid[hi]{pt}{cut}{itp}

\olsection{माएहारा प्रमेय और क्रेग अंतर्वेशन प्रमेय}

विलियम क्रेग का अंतर्वेशन प्रमेय कहता है कि यदि $!A \Entails !B$, तो ऐसा
!!a{sentence}~$!C$ (\emph{अंतर्वेशक}) है कि $!A \Entails !C$ और $!C \Entails
!B$। महत्त्वपूर्ण बात यह है कि अंतर्वेशक~$!C$ को इस प्रकार चुना जा सकता है कि
उसमें आने वाले सभी !!{constant}s और !!{predicate}s,~$!A$ और~$!B$ दोनों में
आएँ। (हम मानते हैं कि कोई !!{function} सम्मिलित नहीं है।) इस प्रतिबंध के बिना
$!A$ और $!B$ तुच्छ रूप से अंतर्वेशक बन सकते थे।

अंतर्वेशन प्रमेय चिरसम्मत प्रथम-कोटि तर्क के लिए सत्य है और इसे खोजा गया
``प्रथम-कोटि तर्क का अंतिम महत्त्वपूर्ण गुण'' कहा गया है। निदर्श सिद्धांत और
स्वचालित तर्कणा में इसके महत्त्वपूर्ण अनुप्रयोग हैं। इसकी मूल प्रेरणा और
अनुप्रयोग विज्ञान-दर्शन में था, क्योंकि इससे उन आदिम पदों का अध्ययन किया जा
सकता है जिनके द्वारा किसी सिद्धांत को अभिगृहीतबद्ध किया जा सकता है या नहीं।
इससे बेथ का परिभाष्यता प्रमेय भी मिलता है, जो कहता है कि यदि कोई सिद्धांत किसी
संकल्पना को परोक्ष रूप से परिभाषित कर सकता है, तो वह उसे प्रत्यक्ष रूप से भी
परिभाषित कर सकता है।

गणितीय तर्क के अनेक परिणामों की तरह अंतर्वेशन प्रमेय को निदर्श सिद्धांत और
प्रमाण सिद्धांत दोनों की विधियों से सिद्ध किया जा सकता है। प्रमाण-सैद्धांतिक
विधियों का लाभ यह है कि वे अंतर्वेशक का अस्तित्व ही नहीं दिखातीं, बल्कि $!A
\Entails !B$ के किसी !!{proof}—जैसे $!A \Sequent !B$ के~\Log{G3c} में किसी
!!{proof}—से अंतर्वेशक की गणना करने वाली कलन-विधि भी देती हैं।

$\Lang L(!A)$ को~$!A$ के !!{constant}s और !!{predicate}s का समुच्चय, और
$\Lang L(\Gamma) = \bigcup \Setabs{\Lang L(!A)}{!A \in \Gamma}$ मानें। इस
पूरे अनुभाग में हम मानते हैं कि संबंधित !!{formula}s में कोई !!{function} नहीं
है।

\begin{thm}[क्रेग अंतर्वेशन प्रमेय]\ollabel{thm:interpolation}
~$!A$ की !!{language} $\Lang L(!A)$ और $!B$ की $\Lang L(!B)$ हो। यदि
$\Log{G3c} \Proves !A \Sequent !B$, तो !!{language}~$\Lang L(!C) \subseteq
\Lang L(!A) \cap \Lang L(!B)$ में ऐसा !!a{formula}~$!C$ है कि $\Log{G3c}
\Proves !A \Sequent !C$ और $\Log{G3c} \Proves !C \Sequent !B$।
\end{thm}

~\Log{G3c} के !!{proof}s की संरचना का उपयोग करने के लिए हम अंतर्वेशन प्रमेय
का सामान्यीकरण सिद्ध करते हैं। इसके लिए अनुक्रमों के पूर्वपक्ष और अनुवर्ती,
दोनों को दो-दो भागों में विभक्त मानें। $\maeh{\Gamma_1}{\Gamma_2} \Sequent
\maeh{\Delta_1}{\Delta_2}$ लिखकर बताते हैं कि $\Gamma_1$ और $\Delta_1$ पहले
भाग तथा $\Gamma_2$ और $\Delta_2$ दूसरे भाग में हैं। तब अंतर्वेशन प्रमेय निम्न
प्रमेय से तुरंत मिलता है।

\begin{lem}[माएहारा प्रमेय]\ollabel{lem:maehara} यदि $\Log{G3c} \Proves
\maeh{\Gamma_1}{\Gamma_2} \Sequent \maeh{\Delta_1}{\Delta_2}$, तो
!!{language}~$\Lang L(!C) \subseteq \Lang L(\Gamma_1, \Delta_1) \cap \Lang
L(\Gamma_2, \Delta_2)$ में ऐसा !!a{formula}~$!C$ है कि $\Log{G3c} \Proves
\Gamma_1 \Sequent \Delta_1, !C$ और $\Log{G3c} \Proves !C, \Gamma_2 \Sequent
\Delta_2$।
\end{lem}

\begin{proof}
  मान लें $\pi \Proves \maeh{\Gamma_1}{\Gamma_2} \Sequent
  \maeh{\Delta_1}{\Delta_2}$। हम कथन को $n=\pheight{\pi}$ पर आगमन से सिद्ध
  करते हैं।

  यदि $n=0$, तो $\maeh{\Gamma_1}{\Gamma_2} \Sequent
  \maeh{\Delta_1}{\Delta_2}$ अभिगृहीत है और कोई परमाण्विक !!{formula}~$!A$,
  $\Gamma_1, \Gamma_2$ और $\Delta_1, \Delta_2$ दोनों में आता है। बाईं ओर
  $!A$ के पहले या दूसरे भाग में और दाईं ओर पहले या दूसरे भाग में होने के
  अनुसार चार स्थितियाँ हैं।
  \begin{enumerate}
    \item $!A \in \Gamma_1$ और $!A \in \Delta_1$। हमें ऐसा~$!C$ चाहिए कि
    $\Gamma_1 \Sequent \Delta_1, !C$ तथा $!C, \Gamma_2 \Sequent \Delta_2$
    दोनों !!{provable} हों। $!A$ बाएँ और दाएँ दोनों ओर आने से पहला अनुक्रम
    ~ $!C$ के किसी भी चुनाव पर पहले से अभिगृहीत है। $!C, \Gamma_2 \Sequent
    \Delta_2$ की !!{provable}ता सुनिश्चित करने का एकमात्र विकल्प $!C \ident
    \lfalse$ लेना है।
    \item $!A \in \Gamma_1$ और $!A \in \Delta_2$। तब $\Gamma_1 \Sequent
    \Delta_1, !A$ और $!A, \Gamma_2 \Sequent \Delta_2$ दोनों अभिगृहीत हैं,
    इसलिए $!C \ident !A$ चुन सकते हैं।
    \item $!A \in \Gamma_2$ और $!A \in \Delta_1$। तब $!A, \Gamma_1
    \Sequent \Delta_1$ तथा $\Gamma_2 \Sequent \Delta_2, !A$ अभिगृहीत होते,
    पर उनमें~$!A$ गलत पक्षों पर है। इसलिए $!A$ अंतर्वेशक नहीं है। किंतु
    $\RightR{\lnot}$ और $\LeftR{\lnot}$ से हम $\Gamma_1 \Sequent \Delta_1,
    \lnot!A$ तथा $\lnot !A, \Gamma_2 \Sequent \Delta_2$ को !!{prove} कर सकते
    हैं।
    \item $!A \in \Gamma_2$ और $!A \in \Delta_2$। अभ्यास।
  \end{enumerate}
  वे स्थितियाँ अभ्यास के लिए छोड़ी जाती हैं जहाँ
  $\maeh{\Gamma_1}{\Gamma_2} \Sequent \maeh{\Delta_1}{\Delta_2}$ इसलिए
  अभिगृहीत है कि $\lfalse \in \Gamma_1$ या $\lfalse \in \Gamma_2$।

  यदि $n>0$, तो $\pi$ किसी तार्किक अनुमान पर समाप्त होता है। प्रयुक्त नियम
  और प्रमुख सूत्र के भाग के अनुसार स्थितियों में भेद करना होगा। प्रत्येक
  स्थिति में हम आगमन परिकल्पना का उपयोग कर सकते हैं: प्रमेय उन सभी
  !!{proof}s~$\pi_1$ के लिए सत्य है जिनमें $\pheight{\pi_1} < n$, और चाहे
  ~$\pi_1$ का अंतिम अनुक्रम दो भागों में किसी भी तरह विभाजित हो। कुछ विशेष
  स्थितियों पर विचार करें।
  \begin{enumerate}
    \item $\pi$, प्रमुख सूत्र~$\lnot !A$ वाले $\RightR{\lnot}$ पर समाप्त होता
    है। दो स्थितियाँ हैं: $\lnot !A$ के पहले या दूसरे भाग में होने के अनुसार
    $\pi$ निम्न में से किसी एक पर समाप्त होता है
    \[
    \AxiomC{}
    \RightLabel{$\pi_1$}
    \Deduce$\maeh{!A, \Gamma_1}{\Gamma_2} \fCenter
      \maeh{\Delta_1}{\Delta_2}$
    \RightLabel{\RightR{\lnot}}
    \UnaryInf$\maeh{\Gamma_1}{\Gamma_2} \fCenter
      \maeh{\Delta_1, \lnot !A}{\Delta_2}$
    \DisplayProof
    \]
    अथवा
    \[
    \AxiomC{}
    \RightLabel{$\pi_1$}
    \Deduce$\maeh{\Gamma_1}{!A, \Gamma_2} \fCenter
      \maeh{\Delta_1}{\Delta_2}$
    \RightLabel{\RightR{\lnot}}
    \UnaryInf$\maeh{\Gamma_1}{\Gamma_2} \fCenter
      \maeh{\Delta_1}{\Delta_2, \lnot !A}$
    \DisplayProof
    \]
    ध्यान दें कि यदि प्रमुख !!{formula}~$\lnot !A$ पहले भाग में है, तो हमने
    सक्रिय !!{formula}~$!A$ को भी आधार-वाक्य के पहले भाग में रखा है। यह चुनाव
    हम कर सकते हैं। यदि सक्रिय सूत्र को दूसरे भाग में रखते, तो आगमन परिकल्पना
    फिर भी लागू होती, किंतु तब हम कोई अंतर्वेशक न खोज पाते (अभ्यास: प्रयास
    करें)।

    पहले उस स्थिति पर विचार करें जहाँ $\lnot !A$ पहले भाग में है। आगमन
    परिकल्पना से~$\pi_1$ के अंतिम अनुक्रम का कोई अंतर्वेशक, अर्थात् ऐसा
    !!a{formula}~$!C_1$, है कि निम्न दोनों
    \begin{align*}
    !A, \Gamma_1 & \Sequent \Delta_1, !C_1 &&\text{और} &
    !C_1, \Gamma_2 & \Sequent \Delta_2
    \intertext{सिद्ध-योग्य हैं। हमें ऐसा !!a{formula}~$!C$ चाहिए कि निम्न
    दोनों !!{provable} हों:}
    \Gamma_1 & \Sequent \Delta_1, \lnot !A, !C &&\text{और} &
    !C, \Gamma_2 & \Sequent \Delta_2
    \end{align*}
    स्पष्टतः हम $!C \ident !C_1$ ले सकते हैं: आगमन परिकल्पना पहले ही बताती
    है कि दायाँ अनुक्रम !!{provable} है, और बायाँ अनुक्रम आगमन परिकल्पना से
    मिले अनुक्रम पर~$\RightR{\lnot}$ लगाने से मिलता है। आगमन परिकल्पना यह भी
    बताती है कि $\Lang L(!C_1) \subseteq \Lang L(!A, \Gamma_1, \Delta_1)
    \cap \Lang L(\Gamma_2, \Delta_2)$। चूँकि $\Lang L(\lnot !A) = \Lang
    L(!A)$ और $\Lang L(!C) = \Lang L(!C_1)$, हमें $\Lang L(!C_1) \subseteq
    \Lang L(\Gamma_1, \Delta_1, \lnot !A) \cap \Lang L(\Gamma_2, \Delta_2)$
    मिलता है।

    दूसरी परिस्थिति में स्थिति थोड़ी भिन्न है, क्योंकि $\lnot !A$ दूसरे भाग
    में है। आधार-वाक्य के सक्रिय !!{formula} को भी दूसरे भाग में रखते हैं और
    आगमन परिकल्पना निम्न पर लगाते हैं
    \begin{align*}
    \Gamma_1 & \Sequent \Delta_1, !C_1 &&\text{और} &
    !C_1, !A, \Gamma_2 & \Sequent \Delta_2
    \intertext{सिद्ध-योग्य हैं। यदि हम इनमें से दूसरे पर फिर
    \RightR{\lnot} नियम लगाएँ, तो निम्न के !!{proof}s मिलते हैं}
    \Gamma_1 & \Sequent \Delta_1, !C_1 &&\text{और} &
    !C_1, \Gamma_2 & \Sequent \Delta_2, \lnot !A
    \end{align*}
    यदि $!C_1$ अंतर्वेशक हो तो ठीक इन्हीं अनुक्रमों को !!{provable} होना चाहिए;
    हम फिर $!C \ident !C_1$ लेते हैं। पहले की तरह $\Lang L(!C_1) \subseteq
    \Lang L(\Gamma_1, \Delta_1) \cap \Lang L(\Gamma_2, \Delta_2, \lnot
    !A)$।
    \item $\pi$, प्रमुख !!{formula}~$!A \lif !B$ वाले $\RightR{\lif}$ पर
    समाप्त होता है। $!A \lif !B$ के पहले या दूसरे भाग में होने के अनुसार फिर
    दो स्थितियाँ हैं। !!{proof}~$\pi$ निम्न में से किसी एक पर समाप्त होता है
    \[
    \AxiomC{}
    \RightLabel{$\pi_1$}
    \Deduce$\maeh{!A, \Gamma_1}{\Gamma_2} \fCenter
      \maeh{\Delta_1, !B}{\Delta_2}$
    \RightLabel{\RightR{\lif}}
    \UnaryInf$\maeh{\Gamma_1}{\Gamma_2} \fCenter
      \maeh{\Delta_1, !A \lif !B}{\Delta_2}$
    \DisplayProof
    \]
    अथवा
    \[
    \AxiomC{}
    \RightLabel{$\pi_1$}
    \Deduce$\maeh{\Gamma_1}{!A, \Gamma_2} \fCenter
      \maeh{\Delta_1}{\Delta_2, !B}$
    \RightLabel{\RightR{\lif}}
    \UnaryInf$\maeh{\Gamma_1}{\Gamma_2} \fCenter
      \maeh{\Delta_1}{\Delta_2, !A \lif !B}$
    \DisplayProof
    \]
    हमने फिर आधार-वाक्य के सक्रिय !!{formula}s को उसी भाग में रखा है जिसमें
    निष्कर्ष का प्रमुख !!{formula} है। पहली परिस्थिति में आगमन परिकल्पना लगाकर
    निम्न के !!{proof}s पाते हैं
    \begin{align*}
    !A, \Gamma_1 & \Sequent \Delta_1, !B, !C_1 &&\text{और} & !C_1, \Gamma_2 & \Sequent \Delta_2
    \intertext{बायाँ अनुक्रम \RightR{\lif} का उपयुक्त आधार-वाक्य है। अतः
    निम्न अनुक्रम !!{provable} हैं:}
    \Gamma_1 & \Sequent \Delta_1, !A \lif !B, !C_1 &&\text{और} & !C_1, \Gamma_2 & \Sequent \Delta_2
    \end{align*}
    इसका अर्थ है कि $!C_1$,~$\pi$ के अंतिम अनुक्रम का भी अंतर्वेशक है, और हम
    एक बार फिर $!C \ident !C_1$ ले सकते हैं।
    
    दूसरी परिस्थिति का उपचार उसी प्रकार होता है।
    \item $\pi$, प्रमुख !!{formula}~$!A \lif !B$ वाले $\LeftR{\lif}$ पर
    समाप्त होता है। यदि $!A \lif !B$ पहले भाग में है, तो !!{proof}~$\pi$
    निम्न पर समाप्त होता है
    \[
    \AxiomC{}
    \RightLabel{$\pi_1$}
    \Deduce$\maeh{\Gamma_1}{\Gamma_2} \fCenter
      \maeh{\Delta_1, !A}{\Delta_2}$
    \AxiomC{}
    \RightLabel{$\pi_2$}
    \Deduce$\maeh{!B, \Gamma_1}{\Gamma_2} \fCenter
      \maeh{\Delta_1}{\Delta_2}$
    \BinaryInf$\maeh{!A \lif !B, \Gamma_1}{\Gamma_2} \fCenter
      \maeh{\Delta_1}{\Delta_2}$
    \DisplayProof
    \]
    आगमन परिकल्पना अब हमें चार !!{provable} अनुक्रम देती है: दो बाएँ
    आधार-वाक्य के अंतर्वेशक $!C_1$ के लिए और दो दूसरे आधार-वाक्य के अंतर्वेशक
    $!C_2$ के लिए:
    \begin{align*}
    \Gamma_1 & \Sequent \Delta_1, !A, !C_1 &&\text{(1)} &
    !C_1, \Gamma_2 & \Sequent \Delta_2 && \text{(2)}\\
    !B, \Gamma_1 & \Sequent \Delta_1, !C_2 &&\text{(3)} &
    !C_2, \Gamma_2 & \Sequent \Delta_2 && \text{(4)}
    \end{align*}
    यदि हम दुर्बलीकरण लगाकर (1) की दाईं ओर $!C_2$ और (3) की दाईं ओर $!C_1$
    जोड़ें, तो अनुक्रम (1) और (3),~\LeftR{\lif} के आधार-वाक्य बन सकते हैं:
    \[
    \AxiomC{}
    \Deduce$\Gamma_1 \fCenter \Delta_1, !A, !C_1, !C_2$
    \AxiomC{}
    \Deduce$!B, \Gamma_1 \fCenter \Delta_1, !C_1, !C_2$
    \RightLabel{\LeftR{\lif}}
    \BinaryInf$!A \lif !B, \Gamma_1 \fCenter \Delta_1, !C_1, !C_2$
    \DisplayProof
    \]
    $\RightR{\lor}$ लगाने से अनुक्रम मिलता है
    \[
    !A \lif !B, \Gamma_1 \fCenter \Delta_1, !C_1 \lor !C_2
    \]
    इससे संकेत मिलता है कि इस स्थिति में $!C_1 \lor !C_2$ हमारा अंतर्वेशक हो
    सकता है। और वास्तव में शेष अनुक्रमों (2) और~(4) से दूसरे आवश्यक अनुक्रम
    को !!{prove} किया जा सकता है:
    \[
    \AxiomC{}
    \Deduce$!C_1, \Gamma_2 \fCenter \Delta_2$
    \AxiomC{}
    \Deduce$!C_2, \Gamma_2 \fCenter \Delta_2$
    \RightLabel{\LeftR{\lor}}
    \BinaryInf$!C_1 \lor !C_2, \Gamma_2 \fCenter \Delta_2$
    \DisplayProof
    \]
    यह सत्यापित करना शेष है कि $!C_1 \lor !C_2$ सही भाषा में है। आगमन
    परिकल्पना देती है कि
    \begin{align*}
    \Lang L(!C_1) & \subseteq
    \Lang L(\Gamma_1, \Delta_1, !A) \cap \Lang L(\Gamma_2, \Delta_2) &\text{और}\\
    \Lang L(!C_2) & \subseteq
    \Lang L(!B, \Gamma_1, \Delta_1) \cap \Lang L(\Gamma_2, \Delta_2)
    \intertext{चूँकि}
    \Lang L(!A) & \subseteq \Lang L(!A \lif !B) &\text{और}\\
    \Lang L(!B) & \subseteq \Lang L(!A \lif !B)
    \intertext{हमें दोनों मिलते हैं}
    \Lang L(!C_1) & \subseteq
    \Lang L(\Gamma_1, \Delta_1, !A \lif !B) \cap \Lang L(\Gamma_2, \Delta_2) &\text{और}\\
    \Lang L(!C_2) & \subseteq
    \Lang L(\Gamma_1, \Delta_1, !A \lif !B) \cap \Lang L(\Gamma_2, \Delta_2)
    \intertext{और फलतः $\Lang L (!C_1 \lor !C_2) = {}$}
    \Lang L(!C_1) \cup \Lang L(!C_2) & \subseteq
    \Lang L(\Gamma_1, \Delta_1, !A \lif !B) \cap \Lang L(\Gamma_2, \Delta_2),
    \end{align*}
    जो वांछित है।

    जहाँ $!A \lif !B$ दूसरे भाग में है वहाँ स्थिति भिन्न है, किंतु उसका उपचार
    इसी प्रकार होता है। आगमन परिकल्पना फिर चार !!{provable} अनुक्रम देती है:
    \begin{align*}
    \Gamma_1 & \Sequent \Delta_1, !C_1 &&\text{(1)} &
    !C_1, \Gamma_2 & \Sequent \Delta_2, !A && \text{(2)}\\
    \Gamma_1 & \Sequent \Delta_1, !C_2 &&\text{(3)} &
    !C_2, !B, \Gamma_2 & \Sequent \Delta_2 && \text{(4)}
    \end{align*}
    अब अनुक्रम (2) और (4)—कुछ दुर्बल किए गए !!{formula}s सहित—~\LeftR{\lif}
    के आधार-वाक्य बनते हैं:
    \[
    \AxiomC{}
    \Deduce$!C_1, !C_2, \Gamma_2 \fCenter \Delta_2, !A$
    \AxiomC{}
    \Deduce$!B, !C_1, !C_2, \Gamma_2 \fCenter \Delta_2$
    \RightLabel{\LeftR{\lif}}
    \BinaryInf$!A \lif !B, !C_1, !C_2, \Gamma_2 \fCenter \Delta_2$
    \DisplayProof
    \]
    हम फिर इसे ऐसे अनुक्रम में बदलना चाहते हैं जिसमें $!C_1$ और $!C_2$ से बना
    एक ही !!{formula} हो, किंतु इस बार वह पूर्वपक्ष में चाहिए (क्योंकि अब हम
    दूसरे भाग से निपट रहे हैं)। $\LeftR{\land}$ लगाने से काम बनता है; हमें
    अंतर्वेशक $!C \ident (!C_1 \land !C_2)$ लेना चाहिए। सौभाग्य से शेष
    अनुक्रमों (1) और (3) से दूसरे भाग के संगत अनुक्रम को !!{prove} कर सकते हैं:
    \[
    \AxiomC{}
    \Deduce$\Gamma_1 \fCenter \Delta_1, !C_1$
    \AxiomC{}
    \Deduce$\Gamma_1 \fCenter \Delta_1, !C_2$
    \RightLabel{\RightR{\land}}
    \BinaryInf$\Gamma_1 \fCenter \Delta_1, !C_1 \land !C_2$
    \DisplayProof
    \]
    पहले की तरह $\Lang L(!C_1 \land !C_2) \subseteq \Lang L(\Gamma_1,
    \Delta_1) \cap \Lang L(\Gamma_2, \Delta_2, !A \lif !B)$।

    \item $\pi$, प्रमुख !!{formula}~$\lforall[x][!A(x)]$ वाले
    $\RightR{\lforall}$ पर समाप्त होता है:
    \[
    \AxiomC{}
    \RightLabel{$\pi_1$}
    \Deduce$\maeh{\Gamma_1}{\Gamma_2} \fCenter
    \maeh{\Delta_1, !A(c)}{\Delta_2}$
    \RightLabel{\RightR{\lforall}}
    \UnaryInf$\maeh{\Gamma_1}{\Gamma_2} \fCenter
      \maeh{\Delta_1, \lforall[x][!A(x)]}{\Delta_2}$
    \DisplayProof
    \]
    अथवा
    \[
    \AxiomC{}
    \RightLabel{$\pi_1$}
    \Deduce$\maeh{\Gamma_1}{\Gamma_2} \fCenter
      \maeh{\Delta_1}{\Delta_2, !A(c)}$
    \RightLabel{\RightR{\lforall}}
    \UnaryInf$\maeh{\Gamma_1}{\Gamma_2} \fCenter
      \maeh{\Delta_1}{\Delta_2, \lforall[x][!A(x)]}$
    \DisplayProof
    \]
    पहली परिस्थिति में आगमन परिकल्पना $\Gamma_1 \Sequent \Delta_1, !A(c),
    !C_1$ और $!C_1, \Gamma_2 \Sequent \Delta_2$ के !!{proof}s देती है। पहले
    पर $\RightR{\lforall}$ लगाकर $\Gamma_1 \Sequent \Delta_1,
    \lforall[x][!A(x)], !C_1$ मिलता है। (आइगेनचर-शर्त पूरी है, क्योंकि वह
    ~ $\pi$ के अंत में $\RightR{\lforall}$ नियम में पहले से पूरी थी।) अतः यदि
    $\Lang L(!C_1) \subseteq \Lang L(\Gamma_1, \Delta_1,
    \lforall[x][!A(x)]) \cap \Lang L(\Gamma_2, \Delta_2)$ की शर्त पूरी हो, तो
    $!C_1$ अंतर्वेशक होगा। आगमन परिकल्पना से $\Lang L(!C_1) \subseteq \Lang
    L(\Gamma_1, \Delta_1, !A(c)) \cap \Lang L(\Gamma_2, \Delta_2)$। इसलिए
    हमें~$c$ की चिंता करनी है—क्या $c$,~$!C_1$ में आ सकता है? नहीं। यदि $c$,
    $!C_1$ में आता, तो $c \in \Lang L(\Gamma_1, \Delta_1, !A(c))$ (जो सत्य
    है) और $c \in \Lang L(\Gamma_2, \Delta_2)$ दोनों होते। किंतु तब $c$,
    ~ $\pi$ के $\RightR{\lforall}$ अनुमान के निष्कर्ष में आता और आइगेनचर-शर्त
    का उल्लंघन होता।
    
    दूसरी परिस्थिति समान है।

    \item $\pi$, प्रमुख !!{formula}~$\lexists[x][!A(x)]$ वाले
    $\RightR{\lexists}$ पर समाप्त होता है। फिर दो परिस्थितियाँ हैं। हम उस
    परिस्थिति पर विचार करते हैं जहाँ $\lexists[x][!A(x)]$ पहले भाग में है और
    दूसरी को अभ्यास के लिए छोड़ते हैं।

    पहली परिस्थिति में $\pi$ निम्न पर समाप्त होता है
    \[
    \AxiomC{}
    \RightLabel{$\pi_1$}
    \Deduce$\maeh{\Gamma_1}{\Gamma_2} \fCenter
      \maeh{\Delta_1, \lexists[x][!A(x)], !A(t)}{\Delta_2}$
    \RightLabel{\RightR{\lexists}}
    \UnaryInf$\maeh{\Gamma_1}{\Gamma_2} \fCenter
      \maeh{\Delta_1, \lexists[x][!A(x)]}{\Delta_2}$
    \DisplayProof
    \]
    आगमन परिकल्पना $\Gamma_1 \Sequent \Delta_1, \lexists[x][!A(x)], !A(t),
    !C_1$ और $!C_1, \Gamma_2 \Sequent \Delta_2$ के !!{proof}s देती है। पहले
    पर $\RightR{\lexists}$ लगाकर $\Gamma_1 \Sequent \Delta_1,
    \lexists[x][!A(x)], !C_1$ मिलता है। अतः यदि $\Lang L(!C_1) \subseteq
    \Lang L(\Gamma_1, \Delta_1, \lexists[x][!A(x)]) \cap \Lang L(\Gamma_2,
    \Delta_2)$ की शर्त पूरी हो, तो $!C_1$ फिर अंतर्वेशक होगा। यहाँ आगमन
    परिकल्पना $\Lang L(!C_1) \subseteq \Lang L(\Gamma_1, \Delta_1,
    \lexists[x][!A(x)], !A(t)) \cap \Lang L(\Gamma_2, \Delta_2)$ देती है।
    याद रखें कि हमने माना है कि कोई !!{function} नहीं है। इसका अर्थ है कि
    पद~$t$ केवल !!a{constant} हो सकता है, अर्थात् $t \ident b$। यदि $b \notin
    \Lang L(!C_1)$, अथवा $b \in \Lang L(\Gamma_1, \Delta_1,
    \lexists[x][!A(x)]) \cap \Lang L(\Gamma_2, \Delta_2)$, तो चिंता की बात
    नहीं: $!C_1$ को अंतर्वेशक ले सकते हैं। किंतु यदि $b \in \Lang L(!C_1)$
    और $b \notin \Lang L(\Gamma_1, \Delta_1, \lexists[x][!A(x)]) \cap \Lang
    L(\Gamma_2, \Delta_2)$ तो? सबसे पहले~$!C_1$ में~$b$ की सभी आवृत्तियाँ
    !!a{variable}~$y$ से बदलकर $!C_1 \ident \Subst{!C_1'}{b}{y}$ लिख सकते
    हैं, जहाँ~$!C_1'$ में~$b$ नहीं है। इसके अतिरिक्त या तो $b \notin \Lang
    L(\Gamma_1, \Delta_1, \lexists[x][!A(x)])$ या $b \notin \Lang
    L(\Gamma_2, \Delta_2)$। क्रमशः $!C \ident \lforall[y][!C_1'(y)]$ या $!C
    \ident \lexists[y][!C_1'(y)]$ ले सकते हैं। पहली स्थिति में हमें मिलता है:
    \[
    \AxiomC{}
    \Deduce$\Gamma_1 \fCenter
      \Delta_1, \lexists[x][!A(x)], !A(b), !C_1'(b)$
    \RightLabel{\RightR{\lexists}}
    \UnaryInf$\Gamma_1 \fCenter \Delta_1, \lexists[x][!A(x)], !C_1'(b)$
    \RightLabel{\RightR{\lforall}}
    \UnaryInf$\Gamma_1 \fCenter
      \Delta_1, \lexists[x][!A(x)], \lforall[y][!C_1'(y)]$
    \DisplayProof
    \]
    और
    \[
    \AxiomC{}
    \Deduce$!C_1'(b), \lforall[y][!C_1'(y)], \Gamma_2 \fCenter \Delta_2$
    \RightLabel{\LeftR{\lforall}}
    \UnaryInf$\lforall[y][!C_1'(y)], \Gamma_2 \fCenter \Delta_2$
    \DisplayProof
    \]
    ऊपर के अनुक्रमों के !!{proof}s आगमन परिकल्पना से हैं (दूसरे में हम
    पूर्वपक्ष में $\lforall[y][!C_1'(y)]$ द्वारा दुर्बलीकरण की ग्राह्यता लगाते
    हैं)। पहले !!{proof} में \RightR{\lforall} की आइगेनचर-शर्त पूरी है, क्योंकि
    $b \notin \Lang L(\Gamma_1, \Delta_1, \lexists[x][!A(x)])$, और $!C_1'$
    को इस प्रकार परिभाषित किया गया है कि उसमें~$b$ न हो।
    
    दूसरी स्थिति में हमें (आगमन परिकल्पना और~$\lexists[y][!C_1'(y)]$ द्वारा
    दुर्बलीकरण की ग्राह्यता से) मिलता है:
    \[
    \AxiomC{}
    \Deduce$\Gamma_1 \fCenter
      \Delta_1, \lexists[x][!A(x)], !A(b), \lexists[y][!C_1'(y)], !C_1'(b)$
    \RightLabel{\RightR{\lexists}}
    \UnaryInf$\Gamma_1 \fCenter
      \Delta_1, \lexists[x][!A(x)], \lexists[y][!C_1'(y)], !C_1'(b)$
    \RightLabel{\RightR{\lexists}}
    \UnaryInf$\Gamma_1 \fCenter
      \Delta_1, \lexists[x][!A(x)], \lexists[y][!C_1'(y)]$
    \DisplayProof
    \]
    और
    \[
    \AxiomC{}
    \Deduce$!C_1'(b), \Gamma_2 \fCenter \Delta_2$
    \RightLabel{\LeftR{\lexists}}
    \UnaryInf$\lexists[y][!C_1'(y)], \Gamma_2 \fCenter \Delta_2$
    \DisplayProof
    \]
    दूसरे !!{proof} में \LeftR{\lexists} की आइगेनचर-शर्त पूरी है, क्योंकि $b
    \notin \Lang L(\Gamma_2, \Delta_2)$, और $!C_1'$ को इस प्रकार परिभाषित
    किया गया है कि उसमें~$b$ न हो।
  \end{enumerate}
  अन्य स्थितियाँ समान हैं; उन्हें अभ्यास के लिए छोड़ते हैं।
\end{proof}

\begin{prob}
मान लें हमारे पास निम्न नियमों वाला संयोजक $\lnand$ हो
\[
\Axiom$\Gamma \fCenter \Delta, !A$
\Axiom$\Gamma \fCenter \Delta, !B$
\RightLabel{\LeftR{\lnand}}
\BinaryInf$!A \lnand !B, \Gamma \fCenter \Delta$
\DisplayProof
\quad
\Axiom$!A, !B, \Gamma \fCenter \Delta$
\RightLabel{\RightR{\lnand}}
\UnaryInf$\Gamma \fCenter \Delta, !A \lnand !B$
\DisplayProof
\]
\olref{lem:maehara} के प्रमाण की उन स्थितियों को पूरा करके दिखाएँ कि माएहारा
प्रमेय अब भी सत्य है जिनमें $\pi$ इन नियमों में से किसी एक पर समाप्त होता है।
\end{prob}

माएहारा प्रमेय स्थापित कर लेने के बाद हम उससे क्रेग अंतर्वेशन प्रमेय सिद्ध कर
सकते हैं।

\begin{proof}[\olref{thm:interpolation} का प्रमाण]
  मान लें $!A \Sequent !B$ !!{provable} है। अंतर्वेशक~$!C$ पाने के लिए
  $\maeh{!A}{\ } \Sequent \maeh{\ }{!B}$ पर \olref{lem:maehara} लगाएँ। हमें
  $\Proves !A \Sequent !C$ और $\Proves !C \Sequent !B$ मिलता है।
\end{proof}

मान लें $\Gamma$, किसी !!{language}~$\Lang L$ में !!{sentence}s का सिद्धांत
है जिसमें कोई $n$-स्थानी {predicate}~$R$ है। हम कहते हैं कि $\Gamma$,~$R$ को
\emph{परोक्ष रूप से परिभाषित करता है}, तभी और केवल तभी जब
\begin{align*}
\Gamma, \Gamma'
& \Entails \lforall[x_1][\dots\lforall[x_n][(R(x_1,\dots,x_n) \liff
R'(x_1, \dots, x_n))]],
\intertext{जहाँ $\Gamma'$,~$\Gamma$ में~$R$ की सभी आवृत्तियों को $R' \notin
\Lang L$ से बदलकर मिला है। $\Gamma$,~$R$ को \emph{प्रत्यक्ष रूप से परिभाषित
करता है}, तभी और केवल तभी जब~$R$ से रहित कोई~$!A(x_1, \dots, x_n)$ ऐसा हो कि}
\Gamma &\Entails
\lforall[x_1][\dots\lforall[x_n][(R(x_1,\dots,x_n) \liff !A(x_1,
\dots, x_n))]].
\end{align*}
स्पष्टतः यदि $\Gamma$, $R$ को प्रत्यक्ष रूप से परिभाषित करता है, तो वह~$R$ को
परोक्ष रूप से भी परिभाषित करता है। विलोम स्पष्ट नहीं है, किंतु फिर भी सत्य है:

\begin{lem}[बेथ परिभाष्यता प्रमेय]
  यदि $\Gamma$,~$R$ को परोक्ष रूप से परिभाषित करता है, तो वह उसे प्रत्यक्ष रूप
  से भी परिभाषित करता है।
\end{lem}

\begin{proof}
  मान लें $\Gamma$,~$R$ को परोक्ष रूप से परिभाषित करता है। $R'$, $\Gamma'$
  ऊपर जैसे हों और $\Lang L' = \Lang L(\Gamma') = \Lang L \setminus \{R\}
  \cup \{R'\}$। परिकल्पना से हमें मिलता है
  \begin{align*}
    \Gamma, \Gamma' & 
    \Sequent \lforall[x_1][\dots\lforall[x_n][(R(x_1,\dots,x_n) 
    \liff R'(x_1, \dots, x_n))]]
  \intertext{$c_1$, \dots,~$c_n$ ऐसे !!{constant}s हों जो~$\Gamma$ में नहीं
  आते। उत्क्रमणीयता प्रमेय से हमें मिलता है}
    \Gamma, \Gamma', R(c_1,\dots,c_n) & \Sequent  R'(c_1, \dots, c_n)
  \intertext{माएहारा प्रमेय को निम्न पर लगाएँ}
    \maeh{\Gamma, R(c_1,\dots,c_n)}{\Gamma'} & \Sequent
      \maeh{\ }{R'(c_1, \dots, c_n)}
  \intertext{जिससे ऐसा~$!A \ident !A(x_1, \dots, x_n)$ मिलता है कि}
    \Gamma, R(c_1,\dots,c_n) & \Sequent !A(c_1, \dots, c_n) \\
    !A(c_1, \dots, c_n), \Gamma' & \Sequent R'(c_1, \dots, c_n)
  \intertext{के !!{proof}s हैं। चूँकि $\Lang L(!A) = \Lang L_1 \cap \Lang
  L_2 = \Lang L(\Gamma) \setminus \{R\}$,~$!A$ में~$R$ नहीं है। दूसरे अनुक्रम
  के !!{proof} में हर जगह $R'$ को~$R$ से बदलकर निम्न का !!a{proof} पाएँ}
    !A(c_1, \dots, c_n), \Gamma & \Sequent R(c_1, \dots, c_n)
  \end{align*}
  निम्न का !!a{proof} पाने के लिए \RightR{\lif} और \RightR{\lforall} लगाएँ
  \[
    \Gamma \Sequent \lforall[x_1][\dots\lforall[x_n][(R(x_1,\dots,x_n) \liff !A(x_1, \dots, x_n))]].
  \]
  इससे दिखता है कि $!A(x_1, \dots, x_n)$,~$\Gamma$ में $!R$ को प्रत्यक्ष रूप
  से परिभाषित करता है।
\end{proof}

अंतर्वेशन के समतुल्य तीसरा परिणाम (जिसे उसकी तरह माएहारा प्रमेय से सिद्ध किया
जा सकता है) यह है: यदि दो संगत सिद्धांत~$\Gamma_1$, $\Gamma_2$ किसी पूर्ण
सिद्धांत $\Gamma$ को समाहित करते हैं, जिसकी भाषा~$\Lang L(\Gamma) \subseteq
\Lang L(\Gamma_1) \cap \Lang L(\Gamma_2)$ है, तो $\Gamma_1 \cup \Gamma_2$
संगत है।

याद रखें कि पूर्ण सिद्धांत वह है जिसकी भाषा के प्रत्येक !!{sentence}~$!A$ के
लिए या तो $\Gamma \Proves !A$ या $\Gamma \Proves \lnot !A$। चूँकि $\Gamma_1$
और~$\Gamma_2$,~$\Gamma$ के संगत विस्तार हैं, $\Gamma$ को संगत होना चाहिए।
फलतः यदि $!A$, भाषा $\Lang L(\Gamma_1) \cap \Lang L(\Gamma_2)$ में कोई
!!a{sentence} है, तो $\Gamma_1 \Entails !A$ और $\Gamma_2 \Entails \lnot !A$
दोनों नहीं हो सकते। इसे देखने के लिए मान लें ऐसा कोई~$!A$ है। यदि $\Gamma_1
\Entails !A$, तो $\Gamma \Entails !A$। अन्यथा~$\Gamma$ की पूर्णता से $\Gamma
\Entails \lnot !A$; और $\Gamma \subseteq \Gamma_1$ होने से $\Gamma_1
\Entails \lnot !A$, जिससे $\Gamma_1$ असंगत हो जाता। अतः $\Gamma \Entails
!A$; और $\Gamma \subseteq \Gamma_2$ होने से $\Gamma_2 \Entails !A$। इसलिए
$\Gamma \Entails/ \lnot !A$, अन्यथा $\Gamma_2$ असंगत होता।

प्रमाण के लिए संयुक्त !!{language}~$\Lang L(\Gamma_1) \cap \Lang L(\Gamma_2)$
में ऐसे !!a{sentence}~$!A$ का न होना ही पर्याप्त है जिसके लिए $\Gamma_1
\Entails !A$ और $\Gamma_2 \Entails \lnot !A$। यहाँ हम माएहारा प्रमेय का
उपयोग करके इसका प्रमाण-सैद्धांतिक प्रमाण देते हैं।

\begin{thm}[रॉबिन्सन संयुक्त संगति प्रमेय]
मान लें $\Gamma_1$, $\Gamma_2$ संगत सिद्धांत हैं, और !!{language} $\Lang
L(\Gamma_1) \cap \Lang L(\Gamma_2)$ में कोई~$!A$ ऐसा नहीं है जिसके लिए
$\Gamma_1 \Entails !A$ और $\Gamma_2 \Entails \lnot !A$। तब $\Gamma_1 \cup
\Gamma_2$ संगत है।
\end{thm}

\begin{proof}
  इसके विपरीत मान लें कि $\Gamma_1 \cup \Gamma_2$ असंगत है। तब अनुक्रम
  $\Gamma_1, \Gamma_2 \Sequent \ $ का कोई !!a{proof} है।
  $\maeh{\Gamma_1}{\Gamma_2} \Sequent \maeh{\ }{\ }$ पर माएहारा प्रमेय लगाकर
  !!{language}~$\Lang L(\Gamma_1) \cap \Lang L(\Gamma_2)$ में ऐसा
  !!a{sentence}~$!A$ पाएँ कि $\Proves \Gamma_1 \Sequent !A$ और $\Proves !A,
  \Gamma_2 \Sequent \quad$। दूसरे से $\Proves \Gamma_2 \Sequent \lnot !A$
  मिलता है। किंतु हमने माना था कि ऐसा कोई !!{sentence} नहीं है।
\end{proof}

ऊपर हमने अप्रकट रूप से माना है कि सिद्धांत $\Gamma$, $\Gamma_1$, $\Gamma_2$
सीमित हैं। संहतता से परिणाम अनंत सिद्धांतों के लिए भी सत्य हैं।

\end{document}
